\section{Exact separation as a hitting-set problem}
\label{sec:hitting-set}

Fix one of the four actions and a finite coordinate dictionary $\Dict$.
For distinct pattern orbits $R,R'$, define the separator mask
\[
  \Sep_{\Dict}(R,R')
  =\{q\in\Dict:q(R)\neq q(R')\}.
\]
The full dictionary separates if and only if every such mask is nonempty.  A
subfamily $F\subseteq\Dict$ separates if and only if
\begin{equation}
  F\cap\Sep_{\Dict}(R,R')\neq\varnothing
  \quad\text{for every unordered pair }\{R,R'\}.
  \label{eq:hitting-set}
\end{equation}
Consequently the minimum fingerprint problem is the transversal-number
problem for the hypergraph of nonempty pair-separator masks.

Each coordinate can take several integer values.  For this reason our primary
language is ``hitting set on unordered orbit-pair separation masks.''  The
classical Test Cover problem provides useful optimization context
\cite{DeBontridderEtAl2003}, but applying that name here requires the explicit
multi-outcome generalization encoded by Equation~\eqref{eq:hitting-set}.

The reported minima count selected coordinates from the declared dictionary.
They are not minima among arbitrary integer-valued combinations of those
coordinates.  Indeed, on a finite pattern set, a sufficiently large
mixed-radix linear combination can encode the entire separating signature in
one integer whenever the full dictionary separates.  The nontrivial
optimization problem here is therefore the sparsest separating
\emph{subfamily} of the fixed coordinate dictionary.

\subsection{Lossless reductions}

Two elementary reductions shrink the constraint family without changing its
transversals.  Duplicate separator masks may initially be represented once.
Second, if $S\subsetneq T$ are separator masks, the constraint ``hit $T$'' is
implied by ``hit $S$'' and $T$ may be deleted.  The resulting inclusion-minimal
antichain is the finite hitting-set instance used for the four-vertex readable
proof.  The five-vertex lower certificate instead uses a selected explicit
pair core; infeasibility of that core at a given budget implies infeasibility
of the complete pair system at the same budget.

\subsection{Upper and lower certificates}

An upper certificate is an explicit list of coordinates.  Its checker
recomputes every value in Equation~\eqref{eq:contained-support-count} and
verifies that all pattern-orbit signature rows are distinct.  Cardinality then
gives the upper bound.

A lower certificate fixes a budget $b$ and proves that no subfamily of size at
most $b$ hits all required masks.  Three locally checkable rules suffice:
\begin{enumerate}
  \item a singleton separator mask forces its sole coordinate;
  \item $k$ pairwise disjoint remaining masks require at least $k$
  coordinates;
  \item if every coordinate hits at most $M$ among $C$ remaining masks, at
  least $\lceil C/M\rceil$ coordinates are required.
\end{enumerate}
If none of these closes a state, branch on one uncovered mask.  Every hitting
set must select at least one of its coordinates, so the children are jointly
exhaustive.  Memoizing identical reduced states turns the tree into a DAG.
Appendix~\ref{app:n5-dag} gives the precise checker contract.

The proof semantics therefore do not depend on an optimizer's status code.
Optimization software was used to discover candidate supports and hard pair
cores for the five-vertex instances.  The accepted lower bounds are the
complete DAGs checked from explicit pattern pairs and recomputed separator
masks.

\subsection{Global versus residual pair systems}

Two pair systems occur in later sections.  The \emph{global} minimum uses all
unordered pairs and every coordinate through the exact degree.  The
\emph{top-degree extension} problem retains every lower-degree coordinate;
only pairs still equal on those coordinates need to be hit by top-degree
coordinates.  Confusing these universes changes the objective; the exact
conditional values are stated separately in Section~\ref{sec:conditional}.
